\section{Large-$\beta_0$ coefficient functions}\label{sect_bubbles}%
%%%%%%%%%%%%%%%%%%%%%%%%%%%%%%%%%%%%%%%%%%%%%%%%%%%%%
%

The leading contributions to the renormalon singularities in the coefficient functions for qITDs are shown in Fig.~\ref{fig:bubblechain},
see Appendix A for technical details. We obtain   
\begin{align}
\label{B[H]}
B[H^\parallel](w)&=\frac{2C_F}{w} \biggl\{\Big[\frac{1+\alpha^2}{1-\alpha}
- \Big(2\alpha\,{}_2F_1(1,2-w,2+w,\alpha)
\notag\\&\quad
 +\bar \alpha (1-w^2)\Big)\alpha^w\,h_0(w,X)\Big]_+
\notag\\ &\quad
+\delta(\bar \alpha)\bigg[\frac{3(w^2-w-1)}{(w+2)(2w-1)}h_0(w,X)-\frac{3}{2}\bigg]\biggr\}
\notag\\[1mm] &\quad
+\widetilde{R}(w),
\\
B[H^\perp](w)&= B[H^\parallel](w)-4C_F \bar \alpha(1+w)\alpha^w\,h_0(w,X)\,,
\end{align}
where $\bar\alpha = 1-\alpha$,
\begin{align}
& h_0(w,X)=X^w \frac{\Gamma(1-w)}{\Gamma(2+w)},
\notag\\
%&G_0(w)=\frac{\Gamma(4+2w)}{6\Gamma(1-w)\Gamma(1+w)\Gamma^2(2+w)},
%\notag\\[4mm]
&X=- \frac{v^2 z^2 \mu^2 e^{5/3}}{4},
\end{align}
and the function $\widetilde{R}(w)$ is defined as the series expansion in terms of another function
\begin{align}
R(w)&=2C_F \biggl\{\Big[\frac{1+\alpha^2}{1-\alpha}\frac{\alpha^wG_0(w)\!-\!1}{w}+\alpha^w \bar \alpha (2\!+\!w)G_0(w)\Big]_+
\notag\\&\quad
+\frac{\delta(\bar \alpha)}{w}\bigg[\frac{3}{2}-\frac{2w+3}{(w+2)(w+1)}G_0(w)\bigg]\biggr\},
\notag\\
G_0(w)&=\frac{\Gamma(4+2w)}{6\Gamma(1-w)\Gamma(1+w)\Gamma^2(2+w)},
\end{align}
such that
\begin{align}
  R(w) = \sum_n w^n R_n\,, \qquad \widetilde{R}(w) = \sum_n \frac{w^n}{(n+1)!} R_n \,.
\end{align}
The Taylor expansion of the Borel transform at $w=0$ gives the perturbative expansion for the coefficient functions in terms of the coupling constant. 
The $\mathcal{O}(\alpha_s)$ term is
\begin{widetext}
\begin{eqnarray}
\label{NLO}
H^{(1)\parallel}(\alpha)&=&2C_F\Bigg[
\bigg(-\mathbf{L}_\mu \frac{1+\alpha^2}{1-\alpha}+\frac{3-8\alpha+3\alpha^2-4\ln\bar \alpha}{1-\alpha}\bigg)_+
+\delta(\bar \alpha)\bigg(\frac{3}{2}\mathbf{L}_\mu+\frac{7}{2}\bigg)\Bigg],
\\
H^{(1)\perp}(\alpha)&=&2C_F\Bigg[
 \bigg(-\mathbf{L}_\mu \frac{1+\alpha^2}{1-\alpha}+\frac{1-4\alpha+\alpha^2-4\ln\bar \alpha}{1-\alpha}\bigg)_+
 +\delta(\bar \alpha)\bigg(\frac{3}{2}\mathbf{L}_\mu+\frac{5}{2}\bigg)\Bigg],
\end{eqnarray}
and the $n_f$-part of the $\mathcal{O}(\alpha_s^2)$ correction reads 
\begin{eqnarray}
\label{NNLO}
H^{(2)\parallel}(\alpha)&=&
-\frac{4C_Fn_f}{3}\biggl\{-\frac{\mathbf{L}_\mu^2}{2}\Big(\frac{1+\alpha^2}{1-\alpha}\Big)_+-\mathbf{L}_\mu \Big[\frac{1+\alpha^2}{1-\alpha}\Big(\ln\alpha+\frac{2}{3}\Big)+4\frac{\alpha+\ln\bar \alpha}{1-\alpha}\Big]_+
+\delta(\bar \alpha)\Big(\frac{3\mathbf{L}_\mu^2}{4}+\frac{19}{4}\mathbf{L}_\mu+\frac{159}{16}\Big)
\\\nonumber && \qquad+\Big[\frac{1+\alpha^2}{1-\alpha}\Big(\frac{41}{18}+\frac{7}{6}\ln\alpha-\frac{\ln^2\alpha}{4}\Big)
-\frac{4}{1-\alpha}\Big(\text{Li}_2(\alpha)+\ln^2\bar \alpha+\frac{8}{3}\ln\bar \alpha+\ln \bar \alpha \ln\alpha\Big)
-\frac{\alpha(13+6\ln \alpha)}{1-\alpha}\Big]_+ \biggr\}.
\\
H^{(2)\perp}(\alpha)&=&H^{(1)\parallel}(\alpha)+\frac{4C_Fn_f}{3}\bar \alpha\Big(2\mathbf{L}_\mu+2\ln\alpha+\frac{10}{3}\Big),
\end{eqnarray}
\end{widetext}
where
\begin{eqnarray}
\mathbf{L}_\mu=\ln\bigg(\frac{-v^2z^2\mu^2}{4e^{-2\gamma_E}}\bigg).
\end{eqnarray}
We have checked that the expressions in \eqref{NLO} reproduce after the Fourier transform \eqref{def:qPDF+pPDF}
the one-loop correction to the qPDFs calculated in \cite{Xiong:2013bka,Stewart:2017tvs}. 

%
\subsection{Singularities of the Borel transform}\label{sect_renormalons}%
%
The structure of singularities of the Borel transform of the coefficient functions is illustrated 
in Fig.~\ref{fig:borelplane}. 
%
\begin{figure}[b]
\centering
 \includegraphics[width=0.45\textwidth]{BVZplots/borelplane}
  %~\includegraphics[width=6.0cm]{BMJplots/fit-t4-qPDF}
\caption{Singularity structure of the Borel transform}
\label{fig:borelplane}
\end{figure}
% 
There are an UV renormalon singularity at $w=1/2$ and a series of IR renormalons at positive integer $w=1,2,\ldots$.
To the accuracy of our calculation (single bubble chain) all singularities are simple poles.
These singularities  obstruct  the Borel integral in \eqref{BorelIntegral} (Borel-non-summable renormalons) and must be 
matched by the nonperturbative corrections. Note, that the $\widetilde{R}(w)$ term in \eqref{B[H]} is an analytic function of $w$. Thus, it is irrelevant for 
the discussion of singularities.

%
\subsubsection{Ultraviolet renormalon at $w= 1/2 $}\label{sect_uvren}%
%
The singularity at $w=1/2$ is generated by the contribution of large momenta in the 
self-energy insertions in the Wilson line and is part of the renormalization factor
\begin{align}
  B[H] &\stackrel{w\to 1/2}{=} \frac{-4C_F}{w-1/2}\sqrt{X}.
\end{align}
This singularity is well-known \cite{Beneke:1994sw} and is in the one-to-one correspondence to the linear UV divergence in the Wilson line's self energy (see also discusion prior to (\ref{n-qITD})). It can be removed by considering  normalized qPDFs, \eqref{n-qITD} or \eqref{hat-qITD},  and will not be considered further in this work.

%
\subsubsection{Infrared renormalons}\label{sect_irren}%
%
The  leading IR renormalon singularity is at $w=1$. 
We obtain
\begin{align}
\label{w=1}
B[H^\parallel](w) &\stackrel{w\to 1}{=} \frac{-4C_F}{1-w}\Big[\alpha + \bar\alpha \ln \bar\alpha\Big]X\,,
\notag\\
B[H^\perp](w) &\stackrel{w\to 1}{=} \frac{-4C_F}{1-w}\Big[\alpha + \bar\alpha \ln \bar\alpha + \alpha\bar\alpha\Big]X\,.
\end{align}  
We remind our notation $\bar\alpha = 1-\alpha$.
These expressions present our main result.

Renormalon singularities at $w=n$ ($n=2,3...$) have a generic form
\begin{align}
B[H](w)&=\frac{2C_F}{n-w}\Big[\alpha^np_{n-1}(\alpha)+\frac{(-1)^n\delta(\bar \alpha)}{n!(n\!-\!2)!n^2(2n\!-\!1)}\Big] X^w,
\end{align}
where $p_n(\alpha)$ is a polynomial of order $n$, e.g. $p_1(\alpha)=(5\alpha-3)/6$, $p_2(\alpha)=(\alpha^2-25\alpha+20)/180$, etc.


%
